% !PRINTED 735
\weberpage{735}

That $T:T_0 = C$ is an integer follows, namely, from formula (13).

If, however, we apply theorem 3. to all elements of the fundamental system $\mathcal{E}_i(r)$, there results a system of integers $e_{i,k}$ and an odd number $e$, such that

\[
\mathcal{E}_i(r)^e = \pm\, \tau_0^{e_{i,0}}\tau_1^{e_{i,1}} \cdots \tau_{\nu'-1}^{e_{i,\nu'-1}}.
\]

If we take the logarithms of this, and then form the regulator $T_0$, there results

\[
e^{\nu'} T_0 = \pm\, \Sigma \pm e_{0,0}\, e_{1,1} \cdots e_{\nu'-1,\nu'-1}\, T,
\]

consequently

\[
e^{\nu'} = \pm\, C\, \Sigma \pm e_{0,0}\, e_{1,1} \cdots e_{\nu'-1,\nu'-1},
\]

and since the left-hand side here is odd, $C$ cannot be even, which was to be proved.

If all conjugate values of the normal unit $\mathcal{E}(r_\beta)$ have one and the same sign, then the same holds of $\mathcal{E}(r_\beta)^e$ as well, and consequently in this case the exponents $e_0, e_1, \ldots, e_{\nu'-1}$ in formula (20) must, by 2., be even numbers. If one then writes formula (20) thus:

\[
\mathcal{E}(r) = \pm\, \tau_0^{e_0}\tau_1^{e_1} \cdots \tau_{\nu'-1}^{e_{\nu'-1}}\, \mathcal{E}(r)^{1-e},
\]

one obtains from it the theorem:

5. \gesperrt{A normal unit of the field $H_m$ which has the same sign as its conjugates is, apart from sign, a square of a normal unit.}

\begin{center}
\S.\,201.

\gesperrt{Fundamental system of units of the field $H_m$.}
\end{center}

It still remains to determine the denominator $D$ in the formula for $B$ [\S.\,197, (6)], which was defined by the equation

\[
E = E'\, D
\]

where $E$ and $E'$ denoted the regulators of the fields $H_m, H_\mu$.

Here, then, it is no longer a question of the normal units, but of the units of the field $H_m$ in general.
